\usepackage[utf8]{inputenc}
\usepackage[main=spanish]{babel}

\usetheme{Madrid}
\useoutertheme[subsection=false]{miniframes}
\usecolortheme{default}

\uselanguage{spanish}
\languagepath{spanish}
\deftranslation[to=spanish]{Lemma}{Lema}
\deftranslation[to=spanish]{Theorem}{Teorema}

\setbeamertemplate{caption}[numbered]


\usepackage{nicefrac}
\usepackage{tikz}

% --- Configuración de enlaces ---

\let\originalhref\href
\renewcommand{\href}[2]{\originalhref{#1}{\color{blue}\ul{#2}}}

% --- Configuración de bibliografía ---
\usepackage[backend=biber]{biblatex} % 'style' puede ser numeric, alphabetic, etc.
\addbibresource{Bibliografia.bib} % Tu archivo de citas (.bib)



\newcommand{\bb}[1]{\mathbb{#1}}
\newcommand{\cc}[1]{\mathcal{#1}}
\newcommand{\ul}[1]{\underline{#1}}
\DeclareMathOperator{\diag}{diag}


\setbeamercolor{bloqueAmarilloCc}{bg=orange!90!yellow, fg=white} % Cabecera
\setbeamercolor{bloqueAmarilloCu}{bg=orange!6!yellow!4!white, fg=black} % Cuerpo

\newenvironment{definitionblock}[1]{%
  \begin{beamerboxesrounded}[upper=bloqueAmarilloCc, lower=bloqueAmarilloCu, shadow=true]{#1}%
}{%
  \end{beamerboxesrounded}%
}


\AtBeginSection[]
{
  \begin{frame}
    \frametitle{Índice}
    \tableofcontents[currentsection]
  \end{frame}
}


% Declarar aplicaciones
% 1. Nombre aplicación
% 2. Dominio
% 3. Codominio
% 4. Variable
% 5. Imagen de la variable
\newcommand{\Func}[5]{
    \begin{equation*}
        \begin{array}{rrll}
            \displaystyle #1:& \displaystyle  #2 & \longrightarrow & \displaystyle  #3\\
               & \displaystyle  #4 & \longmapsto & \displaystyle  #5
        \end{array}
    \end{equation*}
}
\DeclareMathOperator{\tr}{Tr}


\usepackage[makeroom]{cancel}